\documentclass[12pt,a4paper]{exam}

% =====================
% Packages
% =====================
\usepackage[utf8]{inputenc}
\usepackage[T1]{fontenc}
\usepackage[french]{babel}
\usepackage{amsmath,amssymb}
\usepackage{geometry}
\usepackage{multicol}

\geometry{margin=2cm}
\setlength{\parindent}{0pt}

% =====================
% Document
% =====================
\begin{document}

\begin{center}
    {\Large \textbf{Automatismes P1 – QCM - 1STMG}}\\[2mm]
    \textit{Proportions, pourcentages et évolutions}
\end{center}

\textbf{Consigne :} Pour chaque question, une seule réponse est correcte.  
Aucune justification n’est demandée.
    
\vspace{0.6cm}
\hrule
\vspace{0.6cm}

\begin{questions}

% ==================================================
\question \textbf{Exprimer une proportion}

\begin{parts}

\part Dans une classe de 28 élèves, 7 sont absents.  
La proportion d’élèves absents est :
\begin{multicols}{2}
    \begin{choices}
        \choice $\dfrac{7}{21}$
        \vspace{5pt}
        \choice $\dfrac{7}{28}$
        \choice $\dfrac{21}{28}$
        \choice $\dfrac{28}{7}$
    \end{choices}
\end{multicols}

\part Dans une entreprise, 36 salariés sur 120 sont en CDI.  
La proportion de salariés en CDI est :
\begin{multicols}{2}
    \begin{choices}
    \choice $\dfrac{36}{120}$
    \vspace{5pt}
    \choice $\dfrac{120}{36}$
    \choice $\dfrac{84}{120}$
    \choice $\dfrac{36}{84}$
    \end{choices}
\end{multicols}

\end{parts}

% ==================================================
\question \textbf{Calculer un pourcentage à partir d’un effectif}

\begin{parts}

\part Dans un lycée de 900 élèves, 180 sont en STMG.  
Le pourcentage d’élèves en STMG est :
\begin{multicols}{2}
\begin{choices}
\choice $15\%$
\choice $18\%$
\choice $20\%$
\choice $25\%$
\end{choices}
\end{multicols}

\part Une entreprise compte 80 employés, dont 12 sont cadres.  
Le pourcentage de cadres est :
\begin{multicols}{2}
\begin{choices}
\choice $12\%$
\choice $15\%$
\choice $18\%$
\choice $20\%$
\end{choices}
\end{multicols}

\end{parts}

% ==================================================
\question \textbf{Calculer un effectif total à partir d’un pourcentage}

\begin{parts}

\part 24 salariés représentent \textbf{30\%} de l’effectif total.  
L’effectif total est :
\begin{multicols}{2}
\begin{choices}
\choice $72$
\choice $80$
\choice $90$
\choice $120$
\end{choices}
\end{multicols}

\part 15 produits représentent \textbf{25\%} du stock.  
Le nombre total de produits est :
\begin{multicols}{2}
\begin{choices}
\choice $40$
\choice $50$
\choice $60$
\choice $75$
\end{choices}
\end{multicols}

\end{parts}

% ==================================================
\question \textbf{Proportion de proportion}

\begin{parts}

\part Dans un lycée, 40\% des élèves sont en STMG.  
Parmi eux, 25\% sont en première.  
La proportion d’élèves de première STMG dans le lycée est :
\begin{multicols}{2}
\begin{choices}
\choice $10\%$
\choice $15\%$
\choice $25\%$
\choice $65\%$
\end{choices}
\end{multicols}

\part 60\% des employés sont des femmes.  
Parmi elles, 30\% travaillent à temps partiel.  
La proportion de femmes travaillant à temps partiel est :
\begin{multicols}{2}
\begin{choices}
\choice $18\%$
\choice $30\%$
\choice $45\%$
\choice $90\%$
\end{choices}
\end{multicols}

\end{parts}

% ==================================================
\question \textbf{Coefficient multiplicateur et taux d’évolution}

\begin{parts}

\part Une hausse de \textbf{12\%} correspond au coefficient multiplicateur :
\begin{multicols}{2}
\begin{choices}
\choice $0{,}88$
\choice $1{,}12$
\choice $1{,}20$
\choice $2{,}12$
\end{choices}
\end{multicols}

% \part Le coefficient multiplicateur est égal à $0{,}85$.  
% Le taux d’évolution correspondant est :
% \begin{multicols}{2}
% \begin{choices}
% \choice une hausse de $15\%$
% \choice une hausse de $85\%$
% \choice une baisse de $15\%$
% \choice une baisse de $85\%$
% \end{choices}
% \end{multicols}

\part Une baisse de \textbf{18\%} correspond au coefficient multiplicateur :
\begin{multicols}{2}
\begin{choices}
\choice $1{,}18$
\choice $0{,}82$
\choice $0{,}18$
\choice $-0{,}18$
\end{choices}
\end{multicols}

\part Le coefficient multiplicateur est égal à $0{,}92$.  
Le taux d’évolution correspondant est :
\begin{multicols}{2}
\begin{choices}
\choice une hausse de $8\%$
\choice une hausse de $92\%$
\choice une baisse de $8\%$
\choice une baisse de $92\%$
\end{choices}
\end{multicols}

\end{parts}

% ==================================================
\question \textbf{Valeur finale après une évolution}

\begin{parts}

\part Un article coûte 80 €. Il augmente de \textbf{10\%}.  
Son nouveau prix est :
\begin{multicols}{2}
\begin{choices}
\choice 88 €
\choice 90 €
\choice 72 €
\choice 70 €
\end{choices}
\end{multicols}

\part Le chiffre d’affaires d’une entreprise est de 2 000 €.  
Il baisse de \textbf{20\%}.  
La valeur finale est :
\begin{multicols}{2}
\begin{choices}
\choice 1 600 €
\choice 1 800 €
\choice 2 400 €
\choice 2 200 €
\end{choices}
\end{multicols}

\end{parts}

% ==================================================
\question \textbf{Valeur initiale avant une évolution}

\begin{parts}

\part Après une hausse de \textbf{25\%}, un prix est de 125 €.  
Le prix initial était :
\begin{multicols}{2}
\begin{choices}
\choice 100 €
\choice 95 €
\choice 156,25 €
\choice 80 €
\end{choices}
\end{multicols}

\part Après une baisse de \textbf{20\%}, un montant est de 160 €.  
Le montant initial était :
\begin{multicols}{2}
\begin{choices}
\choice 128 €
\choice 180 €
\choice 200 €
\choice 220 €
\end{choices}
\end{multicols}

\end{parts}

\end{questions}

\end{document}
